% ------------------------------------------------------------------------
% Documento Textual de Uma Página: Descrição Geral da Nova Arquitetura
% Enfoque em Engenharia de Software (Compatível com abnTeX2)
% ------------------------------------------------------------------------

\section{Arquitetura de Software e Projeto do Sistema}

Sob a perspectiva da Engenharia de Software, a transição da arquitetura legada do ecossistema \emph{IQArMobi} para a nova arquitetura proposta caracteriza uma refatoração estrutural voltada à eliminação de complexidade acidental, redução de acoplamento e consolidação de serviços. A versão predecessora baseava-se em uma topologia distribuída heterogênea composta por onze containers (incluindo roteamento em \emph{Node-RED}, um serviço \emph{proxy} intermediário em \emph{Node.js/TypeScript}, os componentes \emph{Orion Context Broker}, \emph{IoT-Agent} e \emph{Cygnus} da plataforma FIWARE, além de persistência poliglota redundante em PostgreSQL e MongoDB). Esse encadeamento introduzia múltiplos pontos únicos de falha (\emph{Single Points of Failure}), sobrecarga de serialização HTTP/JSON entre \emph{middlewares} e elevado consumo de memória sem agregar regras de negócio preditivas. Na nova arquitetura, esses intermediários foram substituídos por um \textbf{Monólito Modular Orientado a Eventos} desenvolvido sobre o \emph{framework} \emph{Django} (Python), reduzindo a cadeia de ingestão para uma integração direta via \emph{Webhook} e unificando a persistência em uma única instância relacional \emph{PostgreSQL 16}.

O fluxo de integração adota o estilo arquitetural \textbf{Orientado a Eventos (\emph{Event-Driven Architecture} --- EDA)} combinado ao padrão de transformação de dados na borda de rede. Os quadros de telemetria emitidos pelos nós sensores LoRaWAN (ou injetados pelo ambiente de testes de integração \emph{LWN-Simulator}) trafegam via UDP Semtech até o \emph{ChirpStack Gateway Bridge}, que atua como adaptador de protocolo para o barramento \emph{publish/subscribe} \emph{MQTT (Eclipse Mosquitto)}. No servidor de rede \emph{ChirpStack v4} --- apoiado pelo banco em memória \emph{Redis 7} para gestão de estado de sessão ---, um \emph{Codec JavaScript} encapsulado no perfil do dispositivo deserializa a carga útil bruta e estrutura um \emph{Data Transfer Object (DTO)} em formato JSON. Subsequentemente, mediante um \emph{Webhook HTTP} assíncrono (\texttt{POST /api/chirpstack/webhook/?event=up}), o \emph{ChirpStack} aciona diretamente o controlador do \emph{Django}, que realiza a filtragem idempotente de eventos de controle de rede antes de invocar a camada de domínio.

Internamente, a aplicação \emph{backend} foi estruturada segundo os princípios de \textbf{Separação de Responsabilidades (\emph{Separation of Concerns})} e \textbf{Alta Coesão}, organizando-se em dois contextos delimitados (\emph{Bounded Contexts}): (i) o módulo \texttt{telemetria}, responsável pelos controladores de ingestão, exposição de interfaces analíticas e Mapeamento Objeto-Relacional (\textbf{ORM}) das entidades de domínio (\texttt{SensorNode}, \texttt{AirReading}, \texttt{AgentPrediction} e \texttt{AgentAlert}), utilizando indexação em árvore B (\texttt{db\_index=True}) sob o atributo temporal para otimizar consultas de séries históricas; e (ii) o módulo \texttt{agente}, que encapsula a inteligência de negócio na classe de serviço \texttt{AirQualifyService}. Para atender a requisitos não funcionais estritos de latência, concorrência e disponibilidade, o serviço implementa quatro padrões de projeto:
\begin{enumerate}
    \item \textbf{\emph{Singleton Thread-Safe} com Exclusão Mútua (\texttt{threading.Lock}):} assegura que o grafo computacional do \emph{TensorFlow/Keras} e os transformadores de normalização (\texttt{scikit-learn}) sejam instanciados uma única vez em memória RAM, prevenindo contenção de recursos e sobrecarga de E/S a cada requisição;
    \item \textbf{\emph{Hot-Reloading} Atômico (\emph{Zero-Downtime Update}):} permite a substituição em tempo de execução dos artefatos do modelo neural mediante carregamento prévio em variáveis isoladas e troca atômica de ponteiros em memória sob \emph{lock}, preservando a instância anterior em caso de exceção (\emph{graceful fallback});
    \item \textbf{Inferência \emph{Multi-Task} Unificada:} consolida três saídas analíticas (reconstrução via \emph{Autoencoder}, regressão preditiva de 15 minutos e classificação normativa CONAMA/EPA) em uma única chamada ao grafo computacional, reduzindo em dois terços o custo computacional de inferência;
    \item \textbf{\emph{Notification Throttling} (\emph{Cooldown} de Estado):} controla a emissão de alertas verificando registros não resolvidos nos últimos 15 minutos, garantindo idempotência semântica e prevenindo tempestades de notificações (\emph{alert storms}).
\end{enumerate}

Por fim, os requisitos de \textbf{reprodutibilidade, testabilidade e operação contínua (DevOps)} são atendidos via \textbf{Infraestrutura como Código (IaC)} com \emph{Docker Compose}, empregando verificações ativas de saúde (\emph{healthchecks}) e orquestração determinística de dependências. O sistema implementa \textbf{Provisionamento \emph{Zero-Touch}} no processo de inicialização (\texttt{entrypoint.sh} e \texttt{garantir\_chirpstack.py}), populando de forma automatizada e idempotente esquemas de banco de dados, superusuários, chaves criptográficas OTAA, perfis de rádio \texttt{AU915\_0} e rotas de integração, além de suportar perfis parametrizados de execução para desenvolvimento, simulação ponta a ponta e produção otimizada (\texttt{--noreload}).
